\section{Introduction}
\label{sec:introduction}

\subsection{The Problem: METIS Finds Good but Not Optimal Plans}

METIS multilevel graph partitioning~\citep{b11paper} produces redistricting
plans with low edge-cut in $O(n \log n)$ time.
For a state with $n \approx 2{,}000$ census tracts and $k$ districts,
a single METIS run typically achieves Polsby-Popper compactness at the
70th--80th percentile of the feasible plan distribution.
But METIS is a local optimizer: it terminates at a local minimum of the
edge-cut objective, and different random seeds reach different local minima.

Multi-seed METIS (the \texttt{--search multi} mode) runs $m$ independent
METIS seeds and returns the best plan.
With $m = 50$ seeds, multi-seed METIS achieves 8--12\% Polsby-Popper
improvement over single-seed.
But the improvement saturates quickly: seeds are independent and explore
different basins rather than the neighborhood of any one good plan.

Long \recom{} chains (e.g., the \cs{} algorithm with $T = 600$ steps)
achieve stronger mixing and can escape local minima~\citep{b16paper}.
However, a single long chain explores broadly, spending many steps
in regions of plan space far from the best-found plan.
The chain cannot ``remember'' and return to the best plan it visited.

\subsection{Short-Burst: Concentrated Neighborhood Exploration}

Short-Burst optimization~\citep{cannon2022} addresses this problem directly.
The key idea is to decompose the total step budget into a sequence of
\emph{short bursts}---each burst is a short \recom{} chain that starts
from the \emph{endpoint} of the previous burst.
After each burst, the endpoint is recorded and becomes the starting point
for the next burst.
The final plan is selected from the recorded endpoints by a percentile
criterion.

The restart mechanism means the chain never drifts far from the best-found
plan: each burst explores a small neighborhood, and the best neighborhood
endpoint becomes the new center for the next exploration.
This concentrates probability near low-edge-cut regions of plan space
without sacrificing the Markov property.

Cannon et al.~\citep{cannon2022} introduced Short-Burst in the context
of partisan outcome analysis---using short chains to measure partisan
sensitivity of redistricting plans.
Our use of Short-Burst for minimum-edge-cut compactness optimisation is
a natural extension: we apply the same restart mechanism to the
Polsby-Popper objective rather than a partisan metric.

\subsection{Main Contributions}

This paper makes three contributions:

\begin{enumerate}
\item \textbf{Implementation.}
      We implement \shb{} in \texttt{redist-cli}.
      The search mode is \texttt{--search short-burst};
      the Rust type is \texttt{SeedCompositor::ShortBurst}.
      Per-burst seeds are derived by SHA-256 from the base seed
      and burst index, ensuring reproducibility across runs.

\item \textbf{Empirical comparison.}
      We compare \shb{} to multi-seed METIS, \cs{}, and \be{}
      on NC ($k=14$), WI ($k=8$), and TX ($k=38$) under a fixed budget
      of 1{,}000 total \recom{} steps.
      \shb{} achieves 15--22\% Polsby-Popper improvement, placing it
      between multi-seed METIS and \be{}.

\item \textbf{Sensitivity analysis.}
      We characterise how burst length $\ell$ and burst count $m$ interact.
      Burst length is the critical parameter: short bursts ($\ell = 5$--$20$)
      outperform long bursts ($\ell = 50+$), which drift and approach
      single-chain \recom{} in behaviour.
\end{enumerate}

\subsection{Paper Structure}

Section~\ref{sec:background} reviews \recom{}, \cs{}, \be{}, and the
original Short-Burst paper.
Section~\ref{sec:algorithm} gives the formal algorithm and the chain-seed
derivation.
Section~\ref{sec:comparison} presents the four-way comparison on NC, WI,
and TX.
Section~\ref{sec:sensitivity} gives the burst-length and burst-count
sensitivity analysis.
Section~\ref{sec:conclusion} concludes.
